\documentclass[11pt,a4paper]{article}
\usepackage[utf8]{inputenc}
\usepackage[vietnamese]{babel}
\usepackage{amsmath,amssymb,amsfonts}
\usepackage{geometry}
\geometry{a4paper, top=1.2cm, bottom=1.2cm, left=1.4cm, right=1.4cm}
\usepackage{tikz}
\usetikzlibrary{calc, intersections, angles, arrows.meta, shapes.geometric, positioning}
\usepackage{tcolorbox}
\tcbuselibrary{skins, breakable}
\usepackage{tabularx}
\usepackage{enumitem}
\usepackage{fancyhdr}

\pagestyle{fancy}
\fancyhf{}
\setlength{\headheight}{15pt}
\lhead{\small\textbf{CHUYÊN ĐỀ 9+ | HÌNH HỌC VÀO 10}}
\rhead{\small\textbf{BÀI 05: TRỰC TÂM ẨN \& TAM GIÁC TRỰC TÂM}}
\cfoot{\thepage}

% Custom tcolorbox styles
\tcbset{
    boxrule=0.8pt,
    arc=2.5mm,
    fonttitle=\bfseries\sffamily,
    coltitle=white
}

\newtcolorbox{modelbox}[1]{
    colback=blue!4!white,
    colframe=blue!75!black,
    title=#1
}

\newtcolorbox{formulabox}[1]{
    colback=orange!3!white,
    colframe=orange!85!black,
    title=#1
}

\newtcolorbox{pedagogybox}[1]{
    colback=teal!4!white,
    colframe=teal!80!black,
    title=#1
}

\newtcolorbox{solutionbox}[1]{
    colback=green!4!white,
    colframe=green!60!black,
    title=#1
}

\newtcolorbox{exercisebox}[1]{
    colback=purple!4!white,
    colframe=purple!75!black,
    title=#1
}

\begin{document}

% ==============================================================
% TRANG 1: ĐỀ BÀI VÀ HÌNH VẼ TIKZ TRỰC QUAN
% ==============================================================
\begin{center}
    {\Large\bfseries\color{blue!80!black} TUYỆT PHẨM HÌNH HỌC 9+ VÀO LỚP 10 CHUYÊN \& CÔNG LẬP TOP ĐẦU}\\[2mm]
    {\large\bfseries BÀI HÌNH 05: ĐƯỜNG KÍNH VUÔNG GÓC DÂY CUNG \& KỸ THUẬT TRỰC TÂM ẨN}\\[1mm]
    \textit{(Trích xuất và phát triển từ Đề 05 - Tuyển tập 55 Đề Thi Tuyển Sinh Vào Lớp 10)}
\end{center}

\vspace{1mm}

\begin{modelbox}{I. ĐỀ BÀI TOÁN GỐC (TRÍCH ĐỀ 05)}
Cho đường tròn $(O)$ và dây cung $AB$ không đi qua tâm. Gọi $H$ là trung điểm của $AB$, đường thẳng $HO$ cắt đường tròn $(O)$ tại hai điểm $C$ và $D$ ($C$ thuộc cung nhỏ $AB$). Lấy điểm $M$ thuộc đoạn $AH$ ($M$ khác $A$ và $H$), tia $CM$ cắt đường tròn $(O)$ tại điểm thứ hai là $N$.
\begin{enumerate}[label=\textbf{\alph*)}]
    \item Chứng minh bốn điểm $D, N, M, H$ cùng thuộc một đường tròn.
    \item Gọi $I$ là giao điểm của hai đường thẳng $DN$ và $AB$, $K$ là giao điểm của hai đường thẳng $CI$ và $DM$. Chứng minh $CM \cdot CN = CK \cdot CI$.
    \item \textbf{(Phân hóa 9+):} Chứng minh điểm $M$ là tâm của đường tròn nội tiếp tam giác trực tâm $HKN$.
\end{enumerate}
\end{modelbox}

\vspace{1mm}

\begin{center}
\begin{tikzpicture}[line join = round, line cap=round, >=stealth, font=\footnotesize, scale=1.12]
    % Định nghĩa tọa độ chuẩn xác tuyệt đối
    \coordinate (O) at (0, 0);
    \coordinate (D) at (0, 2.6);
    \coordinate (C) at (0, -2.6);
    \coordinate (H) at (0, -1.5);
    \coordinate (A) at (-2.124, -1.5);
    \coordinate (B) at (2.124, -1.5);
    \coordinate (M) at (-1.1, -1.5);
    \coordinate (N) at (-2.6, 0.0);
    \coordinate (I) at (-4.1, -1.5);
    \coordinate (K) at (-1.301, -2.251);
    
    % Vẽ đường tròn (O)
    \draw[thick, teal] (O) circle (2.6);
    
    % Vẽ tam giác CDI và các cạnh
    \draw[thick] (C) -- (D);
    \draw[thick] (I) -- (B);
    \draw[thick, blue!80!black] (D) -- (I);
    \draw[thick, blue!80!black] (C) -- (I);
    \draw[thick, blue!80!black] (C) -- (N);
    \draw[thick, blue!80!black] (D) -- (K);
    
    % Vẽ tam giác trực tâm HKN (nổi bật màu đỏ cam)
    \draw[thick, red!85!black] (H) -- (K) -- (N) -- cycle;
    \draw[dashed, purple] (K) -- (M);
    \draw[dashed, purple] (N) -- (M);
    
    % Ký hiệu góc vuông tại 3 chân đường cao H, N, K
    \pic[draw, angle radius=2.5mm]{right angle = D--H--A};
    \pic[draw, angle radius=2.5mm]{right angle = D--N--C};
    \pic[draw, angle radius=2.5mm]{right angle = D--K--I};
    
    % Đánh dấu các điểm
    \fill (O) circle (1.5pt) node[right=2pt] {$O$};
    \fill (D) circle (1.5pt) node[above] {$D$};
    \fill (C) circle (1.5pt) node[below] {$C$};
    \fill (H) circle (1.5pt) node[below right=1pt] {$H$};
    \fill (A) circle (1.5pt) node[above left] {$A$};
    \fill (B) circle (1.5pt) node[above right] {$B$};
    \fill[blue!90!black] (M) circle (1.6pt) node[above=2pt] {$M$};
    \fill[red] (N) circle (1.5pt) node[above left] {$N$};
    \fill (I) circle (1.5pt) node[left] {$I$};
    \fill[red] (K) circle (1.5pt) node[right] {$K$};
\end{tikzpicture}
\end{center}

\vspace{1mm}

\begin{pedagogybox}{II. MỤC TIÊU TRỌNG TÂM \& BẢN ĐỒ TƯ DUY 9+}
\begin{itemize}[leftmargin=*, nosep]
    \item \textbf{Cốt lõi hình học:} Khám phá \textbf{"Trực tâm ẩn $M$"} trong tam giác lớn $CDI$. Hai đường cao $IH \perp CD$ và $CN \perp DI$ giao nhau tại $M \implies M$ là trực tâm của $\triangle CDI \implies DM \perp CI$ tại $K$.
    \item \textbf{Đòn bẩy tư duy ngược:} Tam giác $HKN$ chính là tam giác trực tâm của $\triangle CDI$ (tạo bởi ba chân đường cao $H, K, N$). Kế thừa Bổ đề Bài 04, trực tâm của tam giác lớn đồng thời là tâm đường tròn nội tiếp của tam giác trực tâm, do đó $M$ là tâm nội tiếp của $\triangle HKN$!
\end{itemize}
\end{pedagogybox}

\newpage

% ==============================================================
% TRANG 2: SƠ ĐỒ TƯ DUY NGƯỢC HÌNH CÂY THÔNG
% ==============================================================
\begin{center}
    {\Large\bfseries\color{blue!80!black} SƠ ĐỒ TƯ DUY NGƯỢC HÌNH CÂY THÔNG (PINE-TREE FORMULA TREE)}\\[1.5mm]
    {\large\bfseries BẺ KHÓA: CHỨNG MINH $M$ LÀ TÂM ĐƯỜNG TRÒN NỘI TIẾP $\triangle HKN$}\\[1mm]
    \textit{(Quy tắc đọc ngược: Từ bệ đỡ khởi phát ở chân trang $\implies$ Hội tụ tại đỉnh ngọn)}
\end{center}

\vspace{1mm}

\begin{center}
\begin{tikzpicture}[
    x=1cm, y=1cm,
    font=\footnotesize,
    >=Stealth,
    arr_main/.style={-{Stealth[length=3.8mm, width=2.4mm]}, line width=1.6pt, red!85!black},
    arr_sub/.style={-{Stealth[length=3.4mm, width=2.2mm]}, line width=1.3pt, blue!75!black},
    arr_root/.style={-{Stealth[length=3.6mm, width=2.4mm]}, line width=1.4pt, teal!85!black},
    target_node/.style={
        rectangle, draw=red!85!black, fill=red!15, line width=1.4pt, rounded corners=2.5mm,
        inner sep=6pt, align=center, text width=11.5cm
    },
    step_node/.style={
        rectangle, draw=blue!80!black, fill=blue!8, line width=1.1pt, rounded corners=2mm,
        inner sep=5pt, align=center, text width=13.5cm
    },
    aux_node/.style={
        rectangle, draw=orange!90!black, fill=orange!12, line width=1.2pt, rounded corners=2mm,
        inner sep=5pt, align=center, text width=14.8cm
    },
    branch_node/.style={
        rectangle, draw=purple!80!black, fill=purple!8, line width=1.1pt, rounded corners=2mm,
        inner sep=5pt, align=center, text width=7.2cm
    },
    sim_node/.style={
        rectangle, draw=blue!85!black, fill=blue!8, line width=1.1pt, rounded corners=2mm,
        inner sep=5pt, align=center, text width=15.5cm
    },
    base_node/.style={
        rectangle, draw=teal!95!black, fill=teal!20, line width=1.4pt, rounded corners=3mm,
        inner sep=6pt, align=center, text width=16.8cm
    }
]

% Đỉnh cây thông
\node[target_node] (T1) at (0, 0) {
    \textbf{\large\color{red!85!black} [ĐÍCH ĐẾN]:\ $M$ LÀ TÂM NỘI TIẾP TAM GIÁC $HKN$}\\[1mm]
    $\boldsymbol{\iff KM, NM, HM \text{ LÀ CÁC PHÂN GIÁC TRONG CỦA } \triangle HKN}$
};

% Tầng 2
\node[step_node] (T2) at (0, -1.9) {
    $\boldsymbol{\Uparrow} \quad \text{Phân giác trong đỉnh } K:\quad \boldsymbol{\widehat{MKH} = \widehat{MKN} \iff 90^\circ - \widehat{IKH} = 90^\circ - \widehat{CKN} \quad (DK \perp CI)}$
};

% Tầng 3
\node[aux_node] (T3) at (0, -3.8) {
    $\boldsymbol{\Uparrow} \quad \text{Quy về so sánh góc:}\quad \boldsymbol{\widehat{IKH} = \widehat{CKN} \iff \widehat{IMH} = \widehat{CDN}}$\\[1mm]
    \textit{(Hai góc phụ nhau với $90^\circ$ qua góc vuông $\widehat{DKI} = \widehat{DKC} = 90^\circ$)}
};

% Tầng 4: Hai nhánh
\node[branch_node] (T4_left) at (-4.2, -6.1) {
    \textbf{\color{purple!90!black} [NHÁNH 1 -- TỨ GIÁC $HKMI$]:}\\[1mm]
    $\widehat{MHI} = \widehat{MKI} = 90^\circ$\\[0.8mm]
    $\implies HKMI \text{ nội tiếp đường kính } MI$\\[0.8mm]
    $\boldsymbol{\implies \widehat{IKH} = \widehat{IMH}} \quad (\text{cùng chắn } \stackrel{\frown}{IH})$
};

\node[branch_node] (T4_right) at (4.2, -6.1) {
    \textbf{\color{purple!90!black} [NHÁNH 2 -- TỨ GIÁC $DKNC$]:}\\[1mm]
    $\widehat{DKC} = \widehat{DNC} = 90^\circ$\\[0.8mm]
    $\implies DKNC \text{ nội tiếp đường kính } CD$\\[0.8mm]
    $\boldsymbol{\implies \widehat{CKN} = \widehat{CDN}} \quad (\text{cùng chắn } \stackrel{\frown}{CN})$
};

% Tầng 5: Tam giác đồng dạng qua đối đỉnh
\node[sim_node] (T5) at (0, -8.5) {
    $\boldsymbol{\Uparrow} \quad \text{So sánh qua đối đỉnh:}\quad \boldsymbol{\widehat{IMH} = \widehat{DMN} = 90^\circ - \widehat{MDN} = \widehat{CDN}} \quad \left(\triangle IMH \backsim \triangle CDN\right)$
};

% Bệ đỡ khởi phát: Trực tâm ẩn M
\node[base_node] (T6) at (0, -10.6) {
    \textbf{\color{teal!95!black} [GỐC RỄ KHỞI PHÁT -- TRỰC TÂM ẨN $M$]:}\\[1mm]
    $H \text{ là trung điểm } AB \implies CD \perp AB \text{ tại } H \implies \boldsymbol{IH \perp CD \text{ tại } H} \quad (\text{Đường cao 1})\\[0.8mm]
    $N \in (O) \text{ đường kính } CD \implies \boldsymbol{CN \perp DI \text{ tại } N} \quad (\text{Đường cao 2})\\[0.8mm]
    $\boldsymbol{\implies M = IH \cap CN \text{ LÀ TRỰC TÂM } \triangle CDI \implies DK \perp CI \text{ TẠI } K} \quad (\text{Đường cao 3})$
};

% Mũi tên liên kết
\draw[arr_main] (T2.north) -- (T1.south);
\draw[arr_main] (T3.north) -- (T2.south);
\draw[arr_sub] (T4_left.north) to[out=90, in=210] (T3.south);
\draw[arr_sub] (T4_right.north) to[out=90, in=330] (T3.south);
\draw[arr_sub] (T5.north) to[out=150, in=270] (T4_left.south);
\draw[arr_sub] (T5.north) to[out=30, in=270] (T4_right.south);
\draw[arr_root] (T6.north) -- (T5.south);

\end{tikzpicture}
\end{center}

\vspace{1mm}

\begin{formulabox}{III. CÔNG THỨC VÀ BỔ ĐỀ CHỐT HẠ}
\begin{itemize}[leftmargin=*, nosep]
    \item \textbf{Bổ đề đường kính vuông góc dây cung:} $H$ là trung điểm dây cung $AB \implies CD \perp AB$ tại $H$.
    \item \textbf{Kỹ thuật nhận diện Trực tâm ẩn:} Khi xuất hiện hai góc vuông $IH \perp CD$ và $CN \perp DI$, giao điểm $M$ chính là trực tâm của $\triangle CDI$, kéo theo đường thứ ba $DM \perp CI$ tại $K$.
\end{itemize}
\end{formulabox}

\newpage

% ==============================================================
% TRANG 3: LỜI GIẢI XUÔI CHUẨN MỰC & BAREM ĐIỂM
% ==============================================================
\begin{center}
    {\Large\bfseries\color{blue!80!black} LỜI GIẢI XUÔI CHUẨN MỰC \& BAREM ĐIỂM CHI TIẾT}\\[1.5mm]
    {\large\bfseries BÀI TOÁN HÌNH HỌC 9+ (3,0 ĐIỂM - CHUẨN THI TUYỂN SINH VÀO LỚP 10)}
\end{center}

\vspace{1mm}

\begin{solutionbox}{LỜI GIẢI CHI TIẾT TỪNG BƯỚC}
\begin{enumerate}[label=\textbf{\arabic*.}, leftmargin=*]
    \item \textbf{Chứng minh bốn điểm $D, N, M, H$ cùng thuộc một đường tròn (0,75 điểm)}
    \begin{itemize}[leftmargin=*, nosep]
        \item Vì $CD$ đi qua tâm $O$ nên $CD$ là đường kính của $(O)$. Điểm $N \in (O)$ nên $\widehat{CND} = 90^\circ$ (góc nội tiếp chắn nửa đường tròn) $\implies \widehat{DNM} = 90^\circ$. \hfill \textbf{[0,25đ]}
        \item Do $H$ là trung điểm của dây cung $AB$ không đi qua tâm nên đường kính $CD \perp AB$ tại $H$, suy ra $\widehat{DHM} = 90^\circ$. \hfill \textbf{[0,25đ]}
        \item Xét tứ giác $DNMH$ có: $\widehat{DNM} + \widehat{DHM} = 90^\circ + 90^\circ = 180^\circ$.\\
        Do đó tứ giác $DNMH$ nội tiếp đường tròn đường kính $DM$. Bốn điểm $D, N, M, H$ cùng thuộc đường tròn này. \hfill \textbf{[0,25đ]}
    \end{itemize}

    \vspace{1mm}
    \item \textbf{Chứng minh $M$ là trực tâm $\triangle CDI$ và $CM \cdot CN = CK \cdot CI$ (0,75 điểm)}
    \begin{itemize}[leftmargin=*, nosep]
        \item Xét $\triangle CDI$ có $CD \perp AB$ tại $H \implies IH \perp CD$ tại $H$, do đó $IH$ là đường cao hạ từ đỉnh $I$. Lại có $\widehat{CND} = 90^\circ \implies CN \perp DI$ tại $N$, do đó $CN$ là đường cao hạ từ đỉnh $C$. \hfill \textbf{[0,25đ]}
        \item Hai đường cao $IH$ và $CN$ cắt nhau tại điểm $M$. Suy ra $M$ là trực tâm của $\triangle CDI$.\\
        Do đó đường thẳng qua đỉnh thứ ba $D$ và $M$ là đường cao thứ ba: $DM \perp CI$ tại điểm $K = DM \cap CI \implies \widehat{CKM} = 90^\circ$. \hfill \textbf{[0,25đ]}
        \item Xét hai tam giác vuông $\triangle CMK$ và $\triangle CIN$ có:
        $\widehat{CKM} = \widehat{CNI} = 90^\circ$ và $\widehat{ICN}$ chung.\\
        Suy ra $\triangle CMK \backsim \triangle CIN$ (g-g) $\implies \dfrac{CM}{CI} = \dfrac{CK}{CN} \implies CM \cdot CN = CK \cdot CI$ (đpcm). \hfill \textbf{[0,25đ]}
    \end{itemize}

    \vspace{1mm}
    \item \textbf{Chứng minh $KM$ là tia phân giác trong của góc $\widehat{HKN}$ (0,75 điểm)}
    \begin{itemize}[leftmargin=*, nosep]
        \item Tứ giác $HKMI$ có $\widehat{MHI} = \widehat{MKI} = 90^\circ \implies HKMI$ nội tiếp đường tròn đường kính $MI$.\\
        Suy ra $\widehat{IKH} = \widehat{IMH}$ (hai góc nội tiếp cùng chắn cung $IH$). \hfill \textbf{[0,25đ]}
        \item Tứ giác $DKNC$ có $\widehat{DKC} = \widehat{DNC} = 90^\circ \implies DKNC$ nội tiếp đường tròn đường kính $CD$.\\
        Suy ra $\widehat{CKN} = \widehat{CDN}$ (hai góc nội tiếp cùng chắn cung $CN$). \hfill \textbf{[0,25đ]}
        \item Lại có $\widehat{IMH} = \widehat{DMN}$ (đối đỉnh). Trong tam giác vuông $DMN$ vuông tại $N$, ta có $\widehat{DMN} = 90^\circ - \widehat{MDN} = \widehat{CDN}$. Suy ra $\widehat{IKH} = \widehat{CKN}$.\\
        Vì $DM \perp CI$ tại $K$ nên: $\widehat{MKH} = 90^\circ - \widehat{IKH} = 90^\circ - \widehat{CKN} = \widehat{MKN}$.\\
        Vậy $KM$ là tia phân giác trong của góc $\widehat{HKN}$. \hfill \textbf{[0,25đ]}
    \end{itemize}

    \vspace{1mm}
    \item \textbf{Chứng minh điểm $M$ là tâm nội tiếp tam giác $HKN$ (0,75 điểm)}
    \begin{itemize}[leftmargin=*, nosep]
        \item Chứng minh hoàn toàn tương tự đối với đỉnh $N$, ta có $NM$ là tia phân giác trong của góc $\widehat{HNK}$. \hfill \textbf{[0,25đ]}
        \item Tương tự đối với đỉnh $H$, ta cũng có $HM$ là tia phân giác trong của góc $\widehat{KHN}$. \hfill \textbf{[0,25đ]}
        \item Trong tam giác trực tâm $HKN$, ba đường phân giác trong $KM, NM, HM$ cùng đồng quy tại điểm $M$.\\
        Vậy điểm $M$ chính là tâm của đường tròn nội tiếp tam giác trực tâm $HKN$ (đpcm). \hfill \textbf{[0,25đ]}
    \end{itemize}
\end{enumerate}
\end{solutionbox}

\newpage

% ==============================================================
% TRANG 4: BÌNH LUẬN SƯ PHẠM & BÀI TẬP PHÁT TRIỂN
% ==============================================================
\begin{center}
    {\Large\bfseries\color{blue!80!black} BÌNH LUẬN SƯ PHẠM, BẪY TƯ DUY \& BÀI TẬP PHÁT TRIỂN}\\[1.5mm]
    {\large\bfseries GIẢI MÃ NGHỆ THUẬT DẪN DẮT \& MỞ RỘNG TẦM NHÌN HÌNH HỌC 9+}
\end{center}

\vspace{1mm}

\begin{pedagogybox}{I. NGHỆ THUẬT SƯ PHẠM: TẠI SAO LẠI NHẬN DIỆN ĐƯỢC TAM GIÁC CDI?}
\begin{itemize}[leftmargin=*]
    \item \textbf{Cội nguồn tư duy:} Đề bài cho điểm $M$ nằm tùy ý trên đoạn $AH$, tạo cảm giác bài toán là một bài hình động khó lường. Tuy nhiên, hai đường thẳng $AB$ và $CD$ vuông góc với nhau tại $H$, và điểm $N$ thuộc đường tròn đường kính $CD$ tạo ra góc vuông $\widehat{CND} = 90^\circ$.
    \item \textbf{Kỹ thuật quy về trực tâm:} Hai góc vuông này cùng thuộc về các đỉnh của tam giác $CDI$. Khi nhận ra $M = IH \cap CN$ chính là trực tâm của $\triangle CDI$, điểm $M$ từ một điểm "di động đơn lẻ" lập tức trở thành tâm điểm kết nối toàn bộ hệ thống ba đường cao và ba chân đường cao $H, K, N$!
    \item \textbf{4 câu hỏi dẫn dắt của Giáo viên trên lớp:}
    \begin{enumerate}[label=\arabic*., nosep]
        \item Đường kính $CD$ đi qua trung điểm $H$ của dây cung $AB$ cho ta tính chất vuông góc nào?
        \item Góc $\widehat{CND}$ chắn nửa đường tròn đường kính $CD$ bằng bao nhiêu độ?
        \item Hai đoạn thẳng $IH$ và $CN$ vuông góc với những cạnh nào của tam giác $CDI$? Vậy $M$ đóng vai trò gì?
        \item Ba điểm $H, K, N$ tạo thành tam giác gì của tam giác $CDI$? Áp dụng Bổ đề Bài 04 ta suy ra điều gì về điểm $M$?
    \end{enumerate}
\end{itemize}
\end{pedagogybox}

\vspace{1mm}

\begin{pedagogybox}{II. CÁC BẪY TƯ DUY VÀ SAI LẦM HỌC SINH THƯỜNG GẶP}
\begin{itemize}[leftmargin=*, nosep]
    \item \textbf{Bẫy 1 -- Dùng đại số tính độ dài phức tạp:} Nhiều học sinh ở câu b cố gắng tính độ dài các đoạn $CM, CN, CK, CI$ theo $R$ và khoảng cách $OH$. Cách tiếp cận này dẫn tới bế tắc vì tính toán cồng kềnh. Trong khi đó, việc nhận ra trực tâm $M \implies DM \perp CI$ tại $K$ cho ngay tam giác đồng dạng $\triangle CMK \backsim \triangle CIN$ trong đúng 2 dòng!
    \item \textbf{Bẫy 2 -- Nhầm lẫn thứ tự các điểm trên dây cung $AB$:} Cần lưu ý $M \in AH$, $H$ là trung điểm $AB$ và $I = DN \cap AB$. Tùy thuộc vị trí của $M$, điểm $I$ nằm ngoài đoạn $AB$ về phía $A$.
\end{itemize}
\end{pedagogybox}

\vspace{1mm}

\begin{exercisebox}{III. BÀI TẬP PHÁT TRIỂN ĐỈNH CAO (LUYỆN TẬP 9+)}
\textbf{Bài toán phát triển 1 (Đường tròn Euler 9 điểm):}\\
Chứng minh rằng bốn điểm $H, K, N$ và trung điểm của ba đoạn thẳng $MC, MD, MI$ cùng nằm trên một đường tròn (Đường tròn Euler của $\triangle CDI$). Bán kính của đường tròn này bằng một nửa bán kính đường tròn ngoại tiếp $\triangle CDI$.\\[1mm]
\textit{(Gợi ý: Áp dụng định lý về đường tròn Euler 9 điểm đi qua ba chân đường cao và ba trung điểm của các đoạn thẳng nối trực tâm với ba đỉnh).}\\[2mm]
\textbf{Bài toán phát triển 2 (Hệ thống Trực tâm liên hoàn):}\\
Kéo dài $KM$ cắt đường thẳng $CD$ tại điểm $P$. Chứng minh rằng điểm $P$ chính là trực tâm của tam giác $CNI$.\\[1mm]
\textit{(Gợi ý: Chứng minh $IP \perp CN$ và $NP \perp CI$ bằng cách khai thác tính chất trực tâm kép trong tứ giác toàn phần).}
\end{exercisebox}

\end{document}
